\documentclass[11pt,a4paper]{article}

\def\nyear {2025}
\def\nterm {Spring}
\def\nlecturer {Orr Shalit}
\def\ncourse {Complex Analysis 2}

\makeatletter

% packages
\usepackage{amssymb,amsfonts,amsmath,calc,tikz,pgfplots,geometry,mathtools}
\usepackage{color}   % May be necessary if you want to color links
\usepackage[svgnames]{xcolor}
\usepackage[hidelinks]{hyperref}
\usepackage{forest}
\usepackage{commath} % ew
\usepackage{esvect} % \vv makes vectors
\usepackage{centernot} % for the comparison
\usepackage{esdiff}
\usepackage{esint}
\usepackage{amsthm}
\usepackage{fancyhdr}
\usepackage{bm}
\usepackage{witharrows}
\usepackage{bookmark}
\usepackage{tikz-cd}
\usepackage{quiver}
\usepackage{bbm}
\usepackage{textcomp}
\usepackage{float}
\usepackage{gensymb}
\usepackage{cleveref}
\usepackage{environ}
\usepackage{comment}
\usepackage{cancel}
\usepackage{xparse}

% Inevitable cs
\usepackage{listings}
\usepackage[]{algorithm2e}

% tikz libraries
\usetikzlibrary{positioning}
\usetikzlibrary{matrix}
\usetikzlibrary{arrows}
\usetikzlibrary{bending}
\usetikzlibrary{arrows.meta}
\usetikzlibrary{decorations.markings}
\usetikzlibrary{decorations.pathreplacing}
\usetikzlibrary{decorations.markings}
\usetikzlibrary{patterns}
\usetikzlibrary{backgrounds}

% Page style setup
\pagestyle{fancy}
\fancyhfoffset{0pt}
\renewcommand{\headwidth}{\textwidth}
\geometry{margin=1in}
\pgfplotsset{compat=1.18}
\setlength{\headheight}{14.6pt}
\addtolength{\topmargin}{-1.6pt}
\hypersetup{
    colorlinks=false,
    linktoc=section,
    linkcolor=black,
}

%% maketitle setup
\ifx \nauthor\undefined
  \def\nauthor{Yeheli Fomberg}
\else
\fi

\ifx \nemail\undefined
  \def\nemail{\href{mailto:yp@campus.technion.ac.il}
  {\texttt{<yp@campus.technion.ac.il>}}}
\else
\fi

\ifx \ncoursehead \undefined
  \def\ncoursehead{\ncourse}
\fi

\lhead{\emph{\nouppercase{\leftmark}}}
\ifx \nextra \undefined
  \rhead{
    \ifnum\thepage=1
    \else
      \ncoursehead
    \fi}
\else
  \rhead{
    \ifnum\thepage=1
    \else
      \ncoursehead \ (\nextra)
    \fi}
\fi

\def\notes{notes}
\def\problems{problems}

\ifx \ntype \undefined
  \def\ntype{notes}
\else
\fi

\ifx \ntype \notes
  \let\@real@maketitle\maketitle
  \renewcommand{\maketitle}{\@real@maketitle\begin{center}
  \begin{minipage}[c]{0.9\textwidth}\centering\footnotesize
  These notes are not endorsed by the lecturers.
  I have revised them outside lectures to incorporate supplementary
  explanations, clarifications, and material for fun.
  While I have strived for accuracy, any errors or misinterpretations 
  are most likely mine.
  \end{minipage}\end{center}}
\else
  \ifx \ntype \problems
    \let\@real@maketitle\maketitle
    \renewcommand{\maketitle}{\@real@maketitle\begin{center}
    \begin{minipage}[c]{0.9\textwidth}\centering\footnotesize
    These problems are mostly based on the problem sets I was given
    when taking the course.
    
    While I have strived for accuracy and completeness,
    some of the problems would not have been here without the help of
    friends.
    \end{minipage}\end{center}}
  \else
  \fi
\fi


% theorem environments
\theoremstyle{definition}
\newtheorem{definition}{Definition}[section]
\newtheorem{remark}{Remark}[section]
\newtheorem{example}{Example}[section]
\newtheorem{exercise}{Exercise}[section]
\newtheorem{paradox}{Paradox}[section]
\newtheorem{entry}{Entry}[section]
\newtheorem*{solution}{Solution}
\newtheorem*{question}{Question}
\newtheorem*{answer}{Answer}
\newtheorem*{problem}{Problem}
\theoremstyle{plain}
\newtheorem{theorem}{Theorem}[section]
\newtheorem{proposition}[theorem]{Proposition}
\newtheorem{lemma}[theorem]{Lemma}
\newtheorem{corollary}[theorem]{Corollary}
\newenvironment{proofof}[1]{%
    \renewcommand{\proofname}{Proof of #1}%
    \begin{proof}
}{%
    \end{proof}
}
% tikz customization
\pgfarrowsdeclarecombine{twolatex'}{twolatex'}{latex'}{latex'}{latex'}{latex'}
\tikzset{->/.style = {decoration={markings,
                                  mark=at position 0.999
                                  with {\arrow[scale=2]{latex'}}},
                      postaction={decorate}}}
\tikzset{<-/.style = {decoration={markings,
                                  mark=at position 0.001 with {\arrowreversed[scale=2]{latex'}}},
                      postaction={decorate}}}
\tikzset{-->/.style = {decoration={markings,
                                  mark=at position 0.999
                                  with {\arrow[scale=2]{latex'[sep=0pt -2.5]}}},
                      postaction={decorate}}}
\tikzset{<--/.style = {decoration={markings,
                                  mark=at position 0.001
                                  with {\arrowreversed[scale=2]{latex'[sep=0pt -2.5]}}},
                      postaction={decorate}}}
\tikzset{<->/.style = {decoration={markings,
                                  mark=at position 0.001 with {\arrowreversed[scale=2]{latex'}},
                                  mark=at position 0.999 with {\arrow[scale=2]{latex'}},},
                      postaction={decorate}}}
\tikzset{->-/.style = {decoration={markings,
    mark=at position 0.50+0.225/#1 with {\arrow[scale=2]{latex'}}},
                      postaction={decorate}}}
\tikzset{-<-/.style = {decoration={markings,
                                   mark=at position 0.50-0.225/#1 with {\arrowreversed[scale=2]{latex'}}},
                       postaction={decorate}}}
\tikzset{->>/.style = {decoration={markings,
                                  mark=at position 1 with {\arrow[scale=2]{latex'}}},
                      postaction={decorate}}}
\tikzset{<<-/.style = {decoration={markings,
                                  mark=at position 0 with {\arrowreversed[scale=2]{twolatex'}}},
                      postaction={decorate}}}
\tikzset{<<->>/.style = {decoration={markings,
                                   mark=at position 0 with {\arrowreversed[scale=2]{twolatex'}},
                                   mark=at position 1 with {\arrow[scale=2]{twolatex'}}},
                       postaction={decorate}}}
\tikzset{->>-/.style = {decoration={markings,
                                   mark=at position 0.50+0.450/#1 with {\arrow[scale=2]{twolatex'}}},
                       postaction={decorate}}}
\tikzset{-<<-/.style = {decoration={markings,
                                   mark=at position 0.50-0.450/#1 with {\arrowreversed[scale=2]{twolatex'}}},
                       postaction={decorate}}}

\pgfarrowsdeclare{biggertip}{biggertip}{%
  \setlength{\arrowsize}{1pt}
  \addtolength{\arrowsize}{.1\pgflinewidth}
  \pgfarrowsrightextend{0}
  \pgfarrowsleftextend{-5\arrowsize}
}{%
  \setlength{\arrowsize}{1pt}
  \addtolength{\arrowsize}{.1\pgflinewidth}
  \pgfpathmoveto{\pgfpoint{-5\arrowsize}{4\arrowsize}}
  \pgfpathlineto{\pgfpointorigin}
  \pgfpathlineto{\pgfpoint{-5\arrowsize}{-4\arrowsize}}
  \pgfusepathqstroke
}
\tikzset{
	EdgeStyle/.style = {>=biggertip}
}

\tikzset{circ/.style = {fill, circle, inner sep = 0, minimum size = 3}}
\tikzset{scirc/.style = {fill, circle, inner sep = 0, minimum size = 1.5}}
\tikzset{mstate/.style={circle, draw, black, text=black, minimum width=0.7cm}}

\tikzset{eqpic/.style={baseline={([yshift=-.5ex]current bounding box.center)}}}

\definecolor{mblue}{rgb}{0.2, 0.3, 0.8}
\definecolor{morange}{rgb}{1, 0.5, 0}
\definecolor{mgreen}{rgb}{0, 0.4, 0.2}
\definecolor{mred}{rgb}{0.5, 0, 0}

% topology
\DeclareMathOperator{\dfr}{dfr} % deformation retract
\newcommand{\Cells}{\text{Cells}}
\newcommand{\RP}{\mathbb{RP}}
\newcommand{\CP}{\mathbb{CP}}

% order theory
\newcommand{\join}{\vee}
\newcommand{\meet}{\wedge}

% algebraic geometry
\DeclareMathOperator{\Rad}{Rad} % Radical of modules
\DeclareMathOperator{\Spec}{Spec}
% \renewcommand{\div}{\mathrm{div}} % divisors
\DeclareMathOperator{\Mer}{Mer} % Meromorphic functions
\DeclareMathOperator{\Exc}{Exc} % Exceptional divisor
\DeclareMathOperator{\Td}{Td} % Todd classes

% algebra

\DeclareMathOperator{\trdeg}{tr.deg.}
\DeclareMathOperator{\odd}{odd}
\DeclareMathOperator{\even}{even}
\DeclareMathOperator{\lcm}{lcm}
\DeclareMathOperator{\ord}{ord}
\DeclareMathOperator{\codim}{codim}
\DeclareMathOperator{\Ann}{Ann}
\DeclareMathOperator{\Ext}{Ext}
\DeclareMathOperator{\Sym}{Sym}
\DeclareMathOperator{\Out}{Out}
\DeclareMathOperator{\Aut}{Aut}
\DeclareMathOperator{\End}{End}
\DeclareMathOperator{\Inn}{Inn}
\DeclareMathOperator{\Mat}{Mat}
\DeclareMathOperator{\Fix}{Fix}
\DeclareMathOperator{\std}{std}
\DeclareMathOperator{\sgn}{sgn}
\DeclareMathOperator{\adj}{adj}
\DeclareMathOperator{\id}{id}
\DeclareMathOperator{\op}{op}
\DeclareMathOperator{\tr}{tr}
\DeclareMathOperator{\diag}{diag}
\DeclareMathOperator{\rank}{rank}
\DeclareMathOperator{\rk}{rk}
\DeclareMathOperator{\coker}{coker}
\DeclareMathOperator{\Mult}{Mult} % multiplier algebra
\DeclareMathOperator{\Frac}{Frac} % fraction field
\DeclareMathOperator{\Rat}{Rat}
\DeclareMathOperator{\Gal}{Gal}
\newcommand{\ioplus}{\overset{\text{int}}{\oplus}} % interior direct sum
\newcommand{\GG}{\mathcal G} % Galois correspondence
\newcommand{\FF}{\mathcal F} % Galois correspondence
\newcommand{\cupp}{\smile} % cup product
\newcommand{\capp}{\frown} % cap product
\newcommand{\actson}{\curvearrowright}
\newcommand{\bra}{\left\langle}
\newcommand{\ket}{\right\rangle}
\newcommand{\idealin}{\triangleleft\,}
\newcommand{\normalin}{\triangleleft\,}
\newcommand{\ip}[2]{\left\langle #1, #2 \right\rangle}
\newcommand{\bigslant}[2]
{{\raisebox{.2em}{$#1$}\left/\raisebox{-.2em}{$#2$}\right.}}
\newcommand{\properideal}{%
  \mathrel{\ooalign{$\lneq$\cr\raise.22ex\hbox{$\lhd$}\cr}}}

% categories
\newcommand{\cat}[1]{\mathbf{#1}}
\newcommand{\Ch}{\cat{Ch}} % Chain complexes   
\newcommand{\Set}{\cat{Set}}
\newcommand{\Grp}{\cat{Grp}}
\newcommand{\Ab}{\cat{Ab}}
\newcommand{\Ring}{\cat{Ring}}
\newcommand{\CRing}{\cat{CRing}}
\newcommand{\Top}{\cat{Top}}
\newcommand{\Vect}{\cat{Vect}}
\newcommand{\cmod}{\cat{mod}}

%% matrix groups
\newcommand{\GL}{\mathrm{GL}}
\newcommand{\Or}{\mathrm{O}}
\newcommand{\PGL}{\mathrm{PGL}}
\newcommand{\PSL}{\mathrm{PSL}}
\newcommand{\PSO}{\mathrm{PSO}}
\newcommand{\PSU}{\mathrm{PSU}}
\newcommand{\SL}{\mathrm{SL}}
\newcommand{\msl}{\mathrm{sl}}
\newcommand{\SO}{\mathrm{SO}}
\newcommand{\Spin}{\mathrm{Spin}}
\newcommand{\Sp}{\mathrm{Sp}}
\newcommand{\SU}{\mathrm{SU}}
\newcommand{\U}{\mathrm{U}}
\let\char\relax
\newcommand{\char}{\text{char}}

% analysis
\renewcommand{\Re}{\mathrm{Re}}
\renewcommand{\Im}{\mathrm{Im}}
\NewDocumentCommand{\dx}{o}{
  \IfNoValueTF{#1}
    {\dif x}                  % if no optional argument
    {\dif^{#1} \!x}         % if optional argument
}
\NewDocumentCommand{\dy}{o}{
  \IfNoValueTF{#1}
    {\dif y}                  % if no optional argument
    {\dif^{#1} \!y}         % if optional argument
}
\newcommand{\dt}{\dif t}
\newcommand{\du}{\dif u}
\newcommand{\dv}{\dif v}
\newcommand{\dz}{\dif z}
\newcommand{\ds}{\dif s}
\newcommand{\dw}{\dif w}
\newcommand{\dr}{\dif r}
\newcommand{\dm}{\dif m}
\newcommand{\df}{\dif f}
\newcommand{\dV}{\dif V}
\newcommand{\dS}{\dif S}
\newcommand{\dA}{\dif \mathrm{A}}
\newcommand{\dmu}{\dif \mu}
\newcommand{\dnu}{\dif \nu}
\newcommand{\dtheta}{\dif \theta}
\newcommand{\domega}{\dif \omega}
\newcommand{\dsigma}{\dif \sigma}
\newcommand{\dtau}{\dif \tau}
\newcommand{\dvarphi}{\dif \varphi}
\renewcommand{\dif}{\mathrm{d}}
\renewcommand{\Dif}{\mathrm{D}}
\renewcommand{\triangle}{\Delta}
\newcommand{\ptriangle}{\partial \triangle}
\DeclareMathOperator{\im}{im}
\DeclareMathOperator{\cis}{cis}
\DeclareMathOperator{\rot}{rot}
\DeclareMathOperator{\vspan}{span}
\DeclareMathOperator{\grad}{grad}
\DeclareMathOperator{\curl}{curl}
\DeclareMathOperator{\esssup}{esssup}
\newcommand{\hodge}{{\mathnormal{\star}}}
\DeclareMathOperator{\range}{range}
\DeclareMathOperator{\Hom}{Hom}
\DeclareMathOperator{\Int}{Int}
\DeclareMathOperator{\Conv}{Conv} % convex hull
\newcommand{\ind}{\mathbbm{1}}
\DeclareMathOperator{\Res}{Res}
\DeclareMathOperator{\graph}{graph}
\DeclareMathOperator{\diam}{diam}
\DeclareMathOperator{\supp}{supp}
\DeclareMathOperator{\Vol}{Vol} % Volume
\DeclareMathOperator{\Area}{Area} % Area
\DeclareMathOperator{\area}{area}
\DeclareMathOperator{\Ind}{Ind} % Index
\newcommand{\length}{\mathrm{length}}

% logic
\DeclareMathOperator{\MOD}{MOD}
\DeclareMathOperator{\Theory}{Theory}
\newcommand{\notimplies}{%
  \mathrel{{\ooalign{\hidewidth$\not\phantom{=}$\hidewidth\cr$\implies$}}}}

% nice
\newcommand{\half}{\frac{1}{2}}
\newcommand{\pair}{\del}
\newcommand{\taking}[1]{\xrightarrow{#1}}
\newcommand{\otaking}[1]{\xleftarrow{#1}}
\newcommand{\inv}{^{-1}}
\newcommand{\ot}{\leftarrow}
\newcommand{\ninfty}{-\infty}
\newcommand{\pinfty}{+\infty}
\newcommand{\floor}[1]{\left\lfloor #1 \right\rfloor}
\newcommand{\ceil}[1]{\left\lceil #1 \right\rceil}
\newcommand{\viff}{\Big\Updownarrow} % vertical iff
\newcommand{\bmid}{\,\middle\vert\,} % big mid

% probability
\newcommand{\Prob}{\mathbf{P}}
\renewcommand{\vec}[1]{\boldsymbol{\mathbf{#1}}}
\DeclareMathOperator{\Bin}{Bin}
\DeclareMathOperator{\Geo}{Geo}
\DeclareMathOperator{\Poi}{Poi}
\DeclareMathOperator{\Exp}{Exp}
\DeclareMathOperator{\Var}{Var} % Variance
\DeclareMathOperator{\Cov}{Cov}

% special letters
\newcommand{\N}{\mathbb{N}}
\newcommand{\Z}{\mathbb{Z}}
\newcommand{\Q}{\mathbb{Q}}
\newcommand{\R}{\mathbb{R}}
\newcommand{\C}{\mathbb{C}}
\newcommand{\F}{\mathbb{F}}
\newcommand{\E}{\mathbf{E}}
\newcommand{\D}{\mathbb{D}}
\renewcommand{\H}{\mathbb{H}}
\newcommand{\ps}{\mathcal{P}}
\newcommand{\M}{\mathcal{M}}
\renewcommand{\L}{\mathcal{L}}
\renewcommand{\O}{\mathcal{O}}
\renewcommand{\U}{\mathcal{U}}
\newcommand{\Omicron}{O}
\newcommand{\powerset}{\mathcal{P}}

% text
\newcommand{\st}{\text{ s.t. }}
\newcommand{\tand}{\quad \text{and} \quad}
\newcommand{\tor}{\quad \text{or} \quad}
\newcommand{\stand}{\text{ and }}
\newcommand{\stor}{\text{ or }}
\renewcommand{\tt}[1]{\textnormal{\textbf{(#1).}}} %tt=theorem title GET RID OF

% title format
\title{\textbf{\ncourse}}

\ifx \ntype \notes
  \author{Notes taken by \nauthor\, \nemail \\ Based on lectures by \nlecturer}
  \date{\nterm\ \nyear}
\else
  \ifx \ntype \problems
    \author{These problems were solved by \nauthor \\ \nemail}
  \else
  \fi
\fi

\makeatother


\begin{document}
\maketitle

% Insert cool image here


\newpage
\tableofcontents
\newpage

\section{Introduction}
\begin{definition}[The $H^{\infty}(U)$ class]
  Let $U$ be an open subset of $\C$.
  We denote the class of analytic functions on $U$
  by $H(U)$ and define the following class of functions
  \[
    H^{\infty}(U) :=
    \set{f \in H(U) \colon \norm{f}_{\infty} = \sup_{z \in U}
    \abs{f(z)} < \infty}.
  \]
\end{definition}

\begin{definition}[Schur class]
  We define the following class of functions
  \[
    S := \set{f \in H^{\infty}(\D) \mid \norm{f}_{\infty} \le 1}.
  \]
\end{definition}

\begin{proposition}[Schwartz's--Pick lemma]
  \label{prop:schwartz-pick}
  Let $f \in S$ and let $z_1,z_2 \in S$.
  Then
  \[
    \abs{\frac{f(z_1) - f(z_2)}{1 - f(z_1) \overline{f(z_2)}}} \le
    \abs{\frac{z_1 - z_2}{1 - z_1 \overline{z_2}}}.
  \]
\end{proposition}
\begin{proof}
  To be added.
\end{proof}
\begin{remark}
  \Cref{prop:schwartz-pick} is also known as the invariant form of
  the Schwartz lemma.
\end{remark}

\begin{definition}[Pseudohyperbolic metric]
  Let $z,w \in \D$.
  We define the pseudohyperbolic metric by
  \[
    \rho(z,w) := \abs{\frac{z - w}{1 - \overline{w} z}}.
  \]
\end{definition}
\begin{remark}
  The map $\rho$ is invariant under M\"obius transformations.
  That is, if $\phi$ is a M\"obius transformation so that $\phi(w) = 0$
  then we have $\rho(z,w) = \abs{\phi(z)}$.
\end{remark}

\begin{remark}
  Let $f \colon \D \to \D$.
  Then from the Schwartz's--Pick lemma we have
  \[
    \rho(f(z),f(w)) \le \rho(z,w).
  \]
  This means that $f$ is a contraction.
  Moreover, there is an equality if and only if $f$ is
  an automorphism of the disc.
  That is because if $f$ is an automorphism of the disc then $f$ and
  $f^{-1}$ are both contractions which implies the equality.
\end{remark}

\begin{problem}[Pick's interpolation problem]
  \label{problem:picks-interpolaltion-problem}
  Let $z_1,\dots,z_n \in \D$ and let
  $w_1,\dots,w_n \in \D$.
  Does there exist a function 
  $f \in H(\D)$ so that
  $f(z_i) = w_i$ for all $1 \le i \le n$
  and $\abs{f(\lambda)} \le 1$ for all $\lambda \in \D$?
\end{problem}
\begin{remark}
  In case $n = 1$ we can solve the function immediately by
  using $f = \psi_{w_1} \circ \psi_{z_1}$ which first interchanges between
  $0$ and $z_1$ then by $0$ and $w_1$ so $f(z_1) = w_1$ as wanted.
  
  In case $n = 2$ we can do this if and only if
  $\rho(w_1,w_2) \le \rho(z_1,z_2)$.
  An equivalent statement is that
  \[
    \begin{bmatrix}
      \frac {1-|z_{1}|^{2}}{1-|w_{1}|^{2}} &
      \frac {1-{\overline {z_{1}}}z_{2}}{1-{\overline {w_{1}}}w_{2}}
    \\[5pt]
      \frac {1-{\overline {z_{2}}}z_{1}}{1-{\overline {w_{2}}}w_{1}} &
      \frac {1-|z_{2}|^{2}}{1-|w_{2}|^{2}}
    \end{bmatrix}
    \geq 0.
  \]
  % that is because?
\end{remark}

\begin{theorem}[Pick's theorem]
  \label{thm:pick}
  Let $z_1,\dots,z_n \in \D$ and let
  $w_1,\dots,w_n \in \D$.
  Then there exists $f \in H(\D)$ so that
  $f(z_i) = w_i$ for all $1 \le i \le n$ if and only if
  the \textit{Pick matrix}:
  \[
    \begin{bmatrix}
      \cfrac{1 - w_i \overline{w_j}}{1 - z_i \overline{z_j}}
    \end{bmatrix}_{i,j=1}^{n}
  \]
  is positive-semidefinite.
\end{theorem}
\begin{proof}
  When $n = 1$ the condition is that
  \[
    \cfrac{1 - w \overline{w}}{1 - z \overline{z}} \geq 0
  \]
  which holds for any $z,w \in D$.
  When $n = 2$ the condition is that
  \[
    \begin{bmatrix}
      \cfrac{1 - \abs{w_1}^2}{1 - \abs{z_1}^2} &
      \cfrac{1 - w_1 \overline{w_2}}{1 - z_1 \overline{z_2}}
      \\
      \cfrac{1 - w_2 \overline{w_1}}{1 - z_2 \overline{z_1}} &
      \cfrac{1 - \abs{w_2}^2}{1 - \abs{z_2}^2}
    \end{bmatrix} \geq 0
  \]
  which reduces to
  $\rho(w_1,w_2) \le \rho(z_1,z_2)$.
\end{proof}

\begin{definition}[Hardy space]
  We define the Hardy space $H^2(\D)$ as followsing
  \[
    H^2(\D) = H^2 =
    \set{h(z) = \sum_{n=0}^{\infty} a_n z^n \in H(\D) \st
    \norm{h}_2^2 := \sum \abs{a_n}^2 < \infty}.
  \]
\end{definition}
\begin{remark}
  Given a power series $\sum a_n z^n \in H^2$ we have by
  the Cauchy--Scwarz's inequality
  \[
    \abs{\sum a_n z^n} \le
    \del{\sum \abs{a_n}^2}^{1/2}
    \del{\sum \abs{z}^{2n}}^{1/2}
  \]
  which converges for any $z \in \D$.
\end{remark}

\begin{definition}[Inner product]
  Given two formal sums we define their inner product to be
  \[
    \ip{a}{b} =
    \ip{\sum a_n z^n}{\sum b_n z^n}_{H^2} =
    \norm{\sum a_n \overline{b^n}} \le
    \norm{a}_2 \norm{b}_2
  \]
\end{definition}
% define norm and metric

\begin{remark}
  The space $H^2$ is a complete metric space and an inner
  product space so it is a Hilbert space.
\end{remark}

\begin{definition}[Szego kernel]
  The Szego kernel is a function
  $K \colon \D \times \D \to \C$
  defined by
  \[
    K(z,w) := \frac{1}{1 - z \overline{w}}
  \]
  and the kernel function at $w \in \D$ is
  \[
    k_w(z) = k(z,w) = \frac{1}{1 - z \overline{w}}.
  \]
\end{definition}

\begin{proposition}
  Let $k_w \in H^2$,
  let $w \in \D$
  and let $h \in H^2$.
  Then $h(w) = \ip{h}{k_w}$
  and consequently from the Cauchy--Schwarz inequality we obtain
  $\abs{h(w)} \le \frac{1}{\sqrt{1 - \abs{w}^2}} \norm{h}_2$.
\end{proposition}
\begin{proof}
  We have that
  \[
    k_w(z) = k(z,w) = \frac{1}{1 - z \overline{w}}.
  \]
  to be added
\end{proof}

\begin{definition}[Bounded operator]
  Let $X$, $Y$ be normed spaces
  and let $T \colon X \to Y$ be a linear function.
  Then $T$ is said to be bounded if and only if $\norm{T} < \infty$.
  We denote $B(X,Y)$ the space of bounded operators from $X$ to $Y$.
\end{definition}
\begin{remark}
  The dual space of $X^*$ is $X^* = B(X,\C)$.
\end{remark}

\begin{remark}
  Let $w \in \D$ and define
  $ev_w \colon H^2 \to \C$ by
  $ev_w(h) = h(w) = \ip{h}{k_w}$ or $h \mapsto h(w)$.
  The evaluation map of the Hardy space $ev_w \colon H^2 \to \C$ is bounded.
\end{remark}

\begin{definition}[Reproducing kernel Hilbert space]
  Let $H$ be a linear subspace of the space of functions from $X$ to $\C$.
  Then $H$ is said to be a RKHS if $H$ is a Hilbert space and
  point evaluation $ev_x$ is bounded for all $x \in X$.
\end{definition}

By Riesz theorem for all $\phi \in H^*$ there exists $g \in G$
so that for all $h \in H$ we have $\phi(h) = \ip{h}{g}$.
Then %???

\begin{definition}[Kernel function]
  
\end{definition}

% the kernel is semi positive definite

\begin{definition}[Multiplier algebra]
  Let $H$ be an RKHS on $X$.
  Then the multiplier algerba on $H$ is the space
  \[
    \Mult(H) :=
    \set{f \colon X \to \C \mid h \in H \implies f h \in H}.
  \]
\end{definition}

% M_f

\begin{question}
  Can we compute $\Mult(H^2)$?
\end{question}
\begin{answer}
  We know that all polynomials are in $\Mult(H^2)$.
\end{answer}

\begin{proposition}
  \label{prop:temp}
  Let $H$ be a RKHS on $X$ and let $f \in \Mult(H)$.
  Then for all $x \in X$ we have $M_f^* k_x = \overline{f(x)} k_x$.
  Conversely, if $T \in B(H)$ then
  $T k_x = \lambda_x K_x$ which implies that $T = M_f^*$
  and $f \in \Mult(H)$.
\end{proposition}

\begin{corollary}
  Every multiplier is a bounded function
  \[
    \abs{f(x)} \le \norm{M_f^*} = \norm{M_f}.
  \]
\end{corollary}
\begin{corollary}
  We have that $\Mult(H^2) \subseteq H^{\infty}(\D)$
  since $\norm{f}_{\infty} \le \norm{f}_{\Mult}$.
\end{corollary}

\begin{proofof}{\Cref{prop:temp}}
  Let $h \in H$ and let $f \in \Mult(H)$.
  Then
  \begin{align*}
    f(x) h(x) &= \ip{fh}{k_x} \\
              &= \ip{M_f h}{k_x} \\
              &= \ip{h}{M_f^* k_x}.
  \end{align*}
\end{proofof}

% needs rewriting

\begin{definition}[Multiplication algebra norm]
  We define for $f \in \Mult(H)$
  \[
    \norm{f}_{\Mult} :=
    \norm{M_f} =
    \sup_{h \neq 0} \frac{\norm{M_f h}}{\norm{h}} < \infty
  \]
\end{definition}

\begin{proposition}
  Let $H$ be a RKHS on $X$ and $t \in B(H)$
  Then we have 
  \[
    \forall x \in X \quad \exists \lambda_x \in \C \st T K_x = \lambda_x K_x
    \iff
    \exists f \in M(H) \st T = M^*_{f}
  \]
  and in this case $\lambda_x = \overline{f(x)}$.
\end{proposition}

\begin{definition}
  Let $f \colon \D \to \C$ and let $r \in (0,1)$.
  Then we define $f_r(z) = f(rz)$.
\end{definition}
\begin{remark}
  If $f \in H^2$ then $f_r \in H^2$.
  This is because if $f(z) = \sum_{n=0}^{\infty} a_n z_n$
  then we have $f_r(z) = \sum_{n=0}^{\infty} r^n b_n z_n$. %b_n or a_n.
\end{remark}

\begin{proposition}
  Let $f(z) = \sum_{n=0}^{\infty} a_n z_n$.
  Then
  \[
    \norm{f}_{2}^{2} =
    \sum \abs{a_n}^2 =
    \lim_{r \nearrow 1}
    \frac{1}{2\pi} \int_{0}^{2\pi} \abs{f_r\del{e^{it}}}^2 \dt.
  \]
\end{proposition}
\begin{remark}
  We also have that
  \[
    \norm{f}_{2}^{2} =
    \int_{T} \abs{f}^2 \dm
  \]
  where $m$ is the Lebesgue measure on the circle.
\end{remark}
\begin{proof}
  to be added.  
\end{proof}

\begin{proposition}
  Let $f \in H^{\infty}$ and $h \in H^2$.
  We have
  \[
    \norm{f h}_{2}^{2} =
    \lim_{r \nearrow 1}
    \frac{1}{2\pi}
    \int_{0}^{2\pi} \abs{f_r\del{e^{it}} h_r\del{e^{it}}}^2 \dt \le
    \norm{f}^2_{\infty} \norm{h}_{2}^{2}
  \]
  which implies that $f \in \Mult(H)$.
  We also have $\norm{f}_{\Mult} = \norm{M_f} \le \norm{f}_{\infty}$.
\end{proposition}
\begin{corollary}
  We get that $\Mult(H^2) = H^{\infty}$.
\end{corollary}

\begin{theorem}[Pick's theorem]
  \label{thm:pick-again}
  Let $z_1,\dots,z_n \in \D$ and let
  $w_1,\dots,w_n \in \D$.
  Then there exists $f \in H(\D)$ so that
  $f(z_i) = w_i$ for all $1 \le i \le n$ if and only if
  the \textit{Pick matrix}:
  \[
    \begin{bmatrix}
      \cfrac{1 - w_i \overline{w_j}}{1 - z_i \overline{z_j}}
    \end{bmatrix}_{i,j=1}^{n}
  \]
  is positive-semidefinite.
\end{theorem}
\begin{proof}
  Let $z_1,\dots,z_n \in X$ and $w_1,\dots,w_n \in \C$.
  Let $H$ be an RKHS on $X$ so that $H = H^2$.
  If there exists $f \in \Mult(H)$ so that
  \begin{enumerate}
    \item[(1)] $\norm{f}_{\Mult} \le 1$;
    \item[(2)] $f(z_i) = w_i$
  \end{enumerate}
  then
  \[
    \begin{bmatrix}
      K(z_i,z_j) (1 - w_i \overline{w_j})
    \end{bmatrix} \geq 0.
  \]
  That is because
  \[
    \norm{M_f^*} = \norm{M_f} = \norm{f}_{\Mult} \le 1
  \]
  ...to be added
  The opposite direction holds if we assume that the kernel is of
  the form $K(z,w) = \frac{1}{1 - z \overline w}$.
\end{proof}

\begin{theorem}
  We have that $f \in \Mult(H)$ so that $\norm{f}_{\Mult} \le 1$
\end{theorem}

\begin{definition}[Harmonic function]
  Let $u \in C^2(\Omega)$ where $\Omega$ is a region in $\C = \R^2$.
  Then $u$ is said to be harmonic if $\Delta u = u_{xx} + u_{yy}$.
\end{definition}

\begin{proposition}
  A real $u$ on a simply connected region $\Omega$ is harmonic if and only
  if there exists $f \in f \in H(\Omega)$ so that $u = \Re(f) = \Im(i f)$
\end{proposition}
\begin{proof}
  Let $f = u + iv \in H(\Omega)$.
  From the Cauchy--Riemman equations we get $u_x = v_y$ and $v_x = -u_y$
  which means that $u_{xx} + u_{yy} = v_{yx} - v_{xy} = 0$.

  Next suppose that $u$ is harmonic.
  Define
  \[
    V(x,y) :=
    \int v_x \dx + v_y \dy =
    \int_{\Gamma} -u_y \dx + u_x \dy.
  \]
  Using Green's theorem we get that
  \[
    \int_{\partial \Omega} -u_y \dx + u_x \dy =
    \int_{\Omega} \underbrace{\del{u_{xx} + u_{yy}}}_{=\,0} \dx \dy = 0
  \]
  After integrating $V$ we get that
  \begin{gather*}
    \diffp{V}{x}(x,y) =
    \lim_{h \to 0} \frac{1}{h} \int_{x}^{x+h} -u_{y}(t,y) \dt =
    -u_{y}\vert_{(x,y)}
  \end{gather*}
  which completes the proof.
\end{proof}

\begin{corollary}
  Harmonic functions are $C^{\infty}$.
\end{corollary}
\begin{proof}
  % Suppose $U$ is harmonic in $(1+\epsilon) \D$.
  % We have that
  % \[
  %   u = \Re f \tand
  %   f(z) = \sum_{n=0}^{\infty} a_n z^n.
  % \]
  % We can now write
  % \[
  %   u = \frac{f + \overline f}{2} =
  % \]
\end{proof}

\begin{theorem}
  Let $f \in \L^1(T)$ and let $Hf$ be defined on the closed unit disc
  $\overline \D$ by
  \[
    (Hf)(r e^{i \theta}) =
    \begin{cases}
      f(e^{i \theta}), &r = 1 \\
      P[f](r e^{i \theta}), &0 \le r < 1
    \end{cases}.
  \]
  Then $Hf$ is harmonic in $\D$ and $Hf \in C(\overline D)$
\end{theorem}
\begin{remark}
  This solves the Dirichlet's problem.
\end{remark}
\begin{proof}

\end{proof}












\end{document}
